\subsection{Distance of Structured Combination of Vectors}
\label{sec:distance}
The difficulty in using a tensor-span vector $\bm{r}$ is that the set
$\TensorSpace(t)$ is not an affine or linear space.
As a result, a key distance-preservation lemma relied upon by Blaze
no longer applies.
This lemma asserts that if one vector is far from a linear space, then a
random linear combination of this vector with others remains far from the
space with high probability.

Specifically, Blaze relies on the following result
\cite[Lemma~6.1]{eurocrypt/BrehmCFRRZ25}, which is originally due to 
Ben-Sasson \etal~\cite[Theorem~4.1]{ccc/Ben-SassonKS18}:

\begin{lemma}
\label{lemma:distanceBlaze}
{\normalfont \cite[Theorem~4.1]{ccc/Ben-SassonKS18}}
Let $\CodeSpace \subset \Field^\mu$ be a linear space with relative distance $\delta$,
and let $\bar{\bm{C}} \in \Field^{\mu\times t}$ be a matrix whose columns are vectors in $\Field^\mu$.
Assume that at least one column of $\bar{\bm{C}}$ is more than $\epsilon$-far from $\CodeSpace$.
Then for any $0<\iota<\epsilon$ \suchthat~ $\epsilon' \coloneqq \epsilon-\iota < \delta/3$ and a random tensor-spanned vector $\bm{r} \ChooseRandom \Field^t$:
\[
    \Pr_{\bm{r} \ChooseRandom \Field^t}\!\left[
        \delta_H(\bar{\bm{C}}\bm{r}, \CodeSpace) \le \epsilon-\iota
    \right]
    \le
    \frac{1}{\iota|\Field|}.
\]
\end{lemma}

\iffalse
\begin{theorem}
{\normalfont \cite[Lemma~6.1]{eurocrypt/BrehmCFRRZ25}}
Let $\Code: \Field^k \to \Field^\mu$ be a linear code with relative distance $\delta$.
Suppose there exists an MLE-IOPP for the relation $\Relation_\mathsf{MLE}[\Code]$ that, for any proximity parameter $0 < \epsilon' < \delta/3$, achieves soundness error $\kappa^\mathsf{s}$.
Then, for any \( \iota > 0 \), the protocol obtained by applying this MLE-IOPP to instances of the form \( \bm{m} = \bm{M} \bm{r} \), where \( \bm{r} \ChooseRandom \Field^t \), constitutes a $t$-interleaved MLE-IOPP for the following relation
\(
\Relation_{t\mbox{-}\mathsf{MLE\mbox{-}IOPP}}[\Code^t]
\):
\begin{equation}
\label{eq:interleavedIopp}
\left\{
    \left(
    \begin{aligned}
    & \bm{v}^\top \in \Field^t, \bm{x} \in \Field^{\log_2 k}, \\
    & \bm{C} \in \Field^{n \times t}, \epsilon \in \mathbb{R}\\
    \end{aligned}
    \right)
    :~
    % \left(
    \begin{aligned}
        & \bm{M} \in \Field^{k \times t} \\
    \end{aligned}
    % \right)
    ~\left|~
    \begin{aligned}
    & \bm{v} = \bm{l}^\top \bm{M} \text{ , where } \bm{l} = \tensorspan{\bm{x}} \\
    & \wedge \delta(\bm{C}, \Enc(\bm{M})) \le \epsilon t
    \end{aligned}
    \right.
\right\},
\end{equation}
with proximity parameter $\epsilon = \epsilon' + \iota$ and soundness error at most:
\[
\kappa^\mathsf{s} + \frac{1}{\iota |\Field|}.
\]
\end{theorem}
\fi

We now extend this lemma to the following theorem for a tensor-spanned $\bm{r}$.

\begin{theorem}[Distance of Tensor Inner Product]
\label{theorem:tensordistance}
{\normalfont (Variant of \cite[Theorem~4.1]{ccc/Ben-SassonKS18})}
Let $\CodeSpace \subset \Field^\mu$ be a linear space with relative distance $\delta$, and let $\bar{\bm{C}} \in \Field^{\mu \times t}$ be a matrix whose columns are vectors in $\Field^\mu$. Suppose that at least one column of $\bar{\bm{C}}$ is more than $\epsilon$-far from $\CodeSpace$, where $\epsilon \leq \delta/3$. Then for any $0 < \iota < \epsilon/\log_2 t$ and a random tensor-spanned vector $\bm{r} \ChooseRandom \TensorSpace(t)$, it holds that:
\[
    \Pr_{\bm{r} \ChooseRandom \TensorSpace(t)}\!\left[
        \delta(\bar{\bm{C}}\bm{r}, \CodeSpace) \le \epsilon - \iota \cdot \log_2 t
    \right]
    \le
    \frac{\log_2 t}{\iota \cdot (|\Field|-2)}.
\]
\end{theorem}


The full proof of Theorem~\ref{theorem:tensordistance} is provided in Appendix~\ref{appx:prooftensordistance}.
Below, we briefly sketch the proof strategy.

Theorem~\ref{theorem:tensordistance} relies on the following lemma.
\begin{lemma}
\label{lemma:affinedistance}
{\normalfont (Corollary of \cite[Lemma~4.2]{ccc/Ben-SassonKS18})}
Let $\CodeSpace \subset \Field^\mu$ be a linear subspace over a field $\Field$ with relative distance $\delta$. Suppose $\bm{u}^* \in \Field^\mu$ satisfies $\Delta(\bm{u}^*, \CodeSpace) > \epsilon \mu$, and fix any vector $\bm{u} \in \Field^\mu$. For any $0 < \iota < \epsilon$ such that $\epsilon - \iota < \delta/3$, define a set:
\[
    A = \left\{ \alpha \in \Field \mid \delta(\bm{u}^* + \alpha \bm{u}, \CodeSpace) \le \epsilon - \iota \right\}.
\]
Then $|A| \leq 1/\iota$.
\end{lemma}

Intuitively, Lemma~\ref{lemma:affinedistance} states that if a vector is far from a
linear space, then most points in any affine line passing through it remain far
from the space.

Although $\TensorSpace(t)$ is not an affine space, the key idea is to decompose a
tensor-spanned vector into a sequence of affine combinations and apply
Lemma~\ref{lemma:affinedistance} iteratively.

For example, consider $t=4$ and let
$\bar{\bm{C}} = [\bar{\bm{c}}_1\mid \bm{c}_2 \mid \bm{c}_3\mid \bm{c}_4]$, where
$\bar{\bm{c}}_1$ is more than $\epsilon$-far from $\CodeSpace$.
Then, by Lemma~\ref{lemma:affinedistance}, with high probability:
\[
    \delta((1-z_1)\bar{\bm{c}}_1 + z_1 \bm{c}_3, \CodeSpace)
    = \delta(\bar{\bm{c}}_1 + \frac{z_1}{1-z_1} \cdot \bm{c}_3, \CodeSpace)
    > \epsilon - \iota.
\]
Applying Lemma~\ref{lemma:affinedistance} again yields, with high probability:
\begin{align*}
    & \delta((1-z_1)(1-z_2)\bar{\bm{c}}_1 +(1-z_1) z_2 \bm{c}_2 + z_1 (1-z_2) \bm{c}_3 + z_1z_2 \bm{c}_4, \CodeSpace)
    \\
    = &
    \delta([(1-z_1)\bar{\bm{c}}_1 + z_1 \bm{c}_3] + \frac{z_2}{1-z_2} \cdot [ (1-z_1)\bar{\bm{c}}_2 + z_1 \bm{c}_4], \CodeSpace)
    > \epsilon - 2\iota
\end{align*}
This process extends inductively to larger $t$,
yielding the bound in
Theorem~\ref{theorem:tensordistance}.


\iffalse
\begin{lemma}[Tensor random linear combination escapes a subspace]
\label{lemma:tensorcombspace}
Let $\mathbb U \subsetneq \Field^n$ be a linear subspace.
Let $\bar{\bm U} \in \Field^{n\times t}$ be a matrix.
Assume that at least one column of $\bar{\bm U}$ lies outside $\mathbb U$.
Let $\bm r \in \Field^t$ be sampled from $\TensorSpace[n]$.
Then
\[
\Pr\big[ \bar{\bm U}\bm r \in \mathbb U \big]
\;\le\;
\frac{\log_2 t}{|\Field|}.
\]
\end{lemma}


\begin{proof}
Since $\mathbb U$ is a proper linear subspace of $\Field^n$,
there exists a nonzero linear functional
$f: \Field^n \to \Field$ such that
$\mathbb U = \ker(\ell)$.

Let
\[
\bm u := (f(\bar{\bm U}_1), \ldots, f(\bar{\bm U}_t)) \in \Field^t.
\]
By assumption, at least one column of $\bar{\bm U}$ lies outside $\mathbb U$,
hence $\bm u \neq \bm 0$.

We have:
\[
\bar{\bm U}\bm r \in \mathbb U
\quad\Longleftrightarrow\quad
f(\bar{\bm U}\bm r) = \langle \bm{u}, \bm r \rangle = 0.
\]

The quantity $\langle \bm{u}, \bm r \rangle$ is a nonzero multilinear polynomial
in $\log_2 t$ independent random variables.
By the Schwartz--Zippel lemma,
\[
\Pr[\langle \bm{u}, \bm r \rangle = 0] \le \frac{\log_2 t}{|\Field|}.
\qedhere
\]
\end{proof}
\fi


\subsection{Proof of Theorem~\ref{theorem:tensordistance}}
\label{appx:prooftensordistance}

\begin{proof}
Suppose the $i$th column of $\bar{\bm{C}}$, denoted by $\bar{\bm{c}}_i$, is $\epsilon$-far from $\CodeSpace$.
Let $\tau_j$ denote the $j$-th bit of $(i-1)$ in big-endian format.

Given $\bm{z} = (z_1, \dots, z_{\log_2 t})$, we define:
\begin{align*}
\alpha_j = 
\begin{cases}
\frac{z_j}{1 - z_j}, & \text{if } \tau_j = 0, \\
\frac{1 - z_j}{z_j}, & \text{if } \tau_j = 1,
\end{cases}
\quad z_j \in \Field \setminus \{0,1\},
\end{align*}
which is equivalent to:
\begin{align*}
z_j = 
\begin{cases}
\frac{\alpha_j}{1 + \alpha_j}, & \text{if } \tau_j = 0, \\
\frac{1}{1 + \alpha_j}, & \text{if } \tau_j = 1,
\end{cases}
\quad \alpha_j \in \Field \setminus \{0, -1\}.
\end{align*}

Consequently, each tensor-span vector:
\[
    \bm{r} =
    (1-z_1,z_1) \kronecker(1-z_2,z_2) \kronecker \cdots
    \kronecker (1-z_{\log_2},z_{\log_2}).
\]
corresponds uniquely to a vector $\bm{\xi}(\bm{\alpha})$ defined as:
\[
\resizebox{\linewidth}{!}{$
\begin{aligned}
\bm{\xi}(\bm{\alpha})
:=&
    \frac{\bm{r}}{\prod_{j=1}^{\log_2t}[(1-\tau_j)(1-z_j)+\tau_jz_j]}
=
\bigotimes_{j=1}^{\log_2t}~
    (~(1-\tau_j)+\tau_j\alpha_j~,~\tau_j+(1-\tau_j)\alpha_j~).
\end{aligned}
$}
\]

Thus, each $\bm{\alpha} \in (\Field \setminus \{0, -1\})^{\log_2 t}$ uniquely determines a vector $\bm{r} \in \TensorSpace(t)$ via $\bm{\xi}(\bm{\alpha})$, and satisfies $\delta(\bar{\bm{C}} \bm{r}, \CodeSpace) = \delta(\bar{\bm{C}} \bm{\xi}, \CodeSpace)$. Hence, it suffices to prove:
\begin{equation}
\label{eq:projectrho}
\Pr_{\bm{\alpha}\ChooseRandom (\Field \setminus \{0,-1\})^{\log_2 t}}
\left[
    \delta(\bar{C}\bm{\xi},\CodeSpace) < \epsilon - \iota \cdot \log_2 t
\right]
\le
\frac{\log_2 t}{\iota(|\Field|-2)}.
\end{equation}

In what follows, we prove Eq.~\eqref{eq:projectrho}.
For clarity, and without loss of generality, we assume $i = 1$.
In this case:
\[
    \alpha_j =\frac{z_j}{1-j},~\forall j\in[\log_2t],
    \qquad
    \rho(\bm{\alpha}) = (1,\alpha_1)\kronecker \cdots \kronecker (1,\alpha_{\log_2 t}).
\]


We call $\bm{\alpha}$ \emph{good} if:
\[
    \delta(\bar{\bm{C}} \bm{\xi}(\bm{\alpha}), \CodeSpace) \ge \epsilon - \iota \cdot \log_2 t,
\]
and \emph{bad} otherwise.

We define a sequence of intermediate vectors using prefixes of $\bm{\alpha}$:
\begin{align*}
&
    \bm{\xi}_0(\bm{\alpha}) := (1, 0) \kronecker (1,0) \kronecker (1,0)\kronecker \cdots \kronecker (1, 0),
\\
&
    \bm{\xi}_1(\bm{\alpha}) := (1, \alpha_1) \kronecker (1,0) \kronecker (1,0)\kronecker \cdots \kronecker (1, 0),
\\
&
    \bm{\xi}_2(\bm{\alpha}) := (1, \alpha_1) \kronecker (1,\alpha_2) \kronecker (1,0) \kronecker \cdots \kronecker (1, 0),
\\
&
    \cdots
\\
&
    \bm{\xi}_{\log_2 n}(\bm{\alpha}) := (1, \alpha_1) \kronecker (1,\alpha_2) \kronecker (1,\alpha_3) \kronecker \cdots \kronecker (1, \alpha_{\log_2 n}).
\end{align*}
Note that $\bm{\xi}_{\log_2 n}(\bm{\alpha})= \bm{\xi}(\bm{\alpha})$.

We say $\bm{\xi}_j(\bm{\alpha})$ is \emph{good} if:
\[
    \delta(\bar{\bm{C}} \bm{\xi}_j(\bm{\alpha}), \CodeSpace) \ge \epsilon - \iota \cdot j,
\]
and \emph{bad} otherwise.
Note that $\bm{\xi}_0(\bm{\alpha})$ is always good because:
\[
    \delta(\bar{\bm{C}} \bm{\xi}_0(\bm{\alpha}), \CodeSpace) = \delta(\bar{\bm{c}}_1, \CodeSpace) \ge \epsilon.
\]


We now analyze the probability that $\bm{\xi}(\bm{\alpha})$ is bad:
{\allowdisplaybreaks
\begin{align}
\label{eq:serialprobability}
\nonumber
& \Pr_{\bm{\alpha}}[~\bm{\xi}(\bm{\alpha}) \text{ is bad}~]
\\ \nonumber
= &
    \Pr_{\bm{\alpha}}[~\bm{\xi}_{\log_2 n}(\bm{\alpha}) \text{ is bad} \mid \bm{\xi}_{\log_2 n -1}(\bm{\alpha}) \text{ is good} ~] \cdot 
    \Pr_{\bm{\alpha}}[~\bm{\xi}_{\log_2 n -1}(\bm{\alpha}) \text{ is good}~]
\\ \nonumber
&
    +    \Pr_{\bm{\alpha}}[~\bm{\xi}_{\log_2 n}(\bm{\alpha}) \text{ is bad} \mid \bm{\xi}_{\log_2 n -1}(\bm{\alpha}) \text{ is bad} ~] \cdot 
    \Pr_{\bm{\alpha}}[~\bm{\xi}_{\log_2 n -1}(\bm{\alpha}) \text{ is bad}~]
\\ \nonumber
\le &
    \Pr_{\bm{\alpha}}[~\bm{\xi}_{\log_2 n}(\bm{\alpha}) \text{ is bad} \mid \bm{\xi}_{\log_2 n -1}(\bm{\alpha}) \text{ is good} ~]
\\ \nonumber
&
    +
    \Pr_{\bm{\alpha}}[~\bm{\xi}_{\log_2 n -1}(\bm{\alpha}) \text{ is bad}~]
\\ \nonumber
\le &
    \Pr_{\bm{\alpha}}[~\bm{\xi}_{\log_2 n}(\bm{\alpha}) \text{ is bad} \mid \bm{\xi}_{\log_2 n -1}(\bm{\alpha}) \text{ is good} ~]
\\ \nonumber
&
    + \Pr_{\bm{\alpha}}[~\bm{\xi}_{\log_2 n - 1}(\bm{\alpha}) \text{ is bad} \mid \bm{\xi}_{\log_2 n -2}(\bm{\alpha}) \text{ is good} ~]
\\ \nonumber
&
    +
    \Pr_{\bm{\alpha}}[~\bm{\xi}_{\log_2 n - 2}(\bm{\alpha}) \text{ is bad}~]
\\ \nonumber
\le &
    \cdots
\\ \nonumber
\le &
    \Pr_{\bm{\alpha}}[~\bm{\xi}_{\log_2 n}(\bm{\alpha}) \text{ is bad} \mid \bm{\xi}_{\log_2 n -1}(\bm{\alpha}) \text{ is good} ~]
\\ \nonumber
&
+ \cdots
\\ \nonumber
&
    + \Pr_{\bm{\alpha}}[~\bm{\xi}_{1}(\bm{\alpha}) \text{ is bad} \mid \bm{\xi}_{0}(\bm{\alpha}) \text{ is good} ~]
\\ \nonumber
&
    +
    \Pr_{\bm{\alpha}}[~\bm{\xi}_{0}(\bm{\alpha}) \text{ is bad}~]
\\
=&
\sum_{j=1}^{\log_2 n}
\Pr_{\bm{\alpha}}[~\bm{\xi}_{j}(\bm{\alpha}) \text{ is bad} \mid \bm{\xi}_{j-1}(\bm{\alpha}) \text{ is good} ~].
\end{align}
}

Thus, it suffices to show that for all $j \in [\log_2 t]$:
\begin{equation}
\label{eq:conditionalprobability}
    \Pr_{\bm{\alpha}}[~\bm{\xi}_{j}(\bm{\alpha}) \text{ is bad} \mid \bm{\xi}_{j-1}(\bm{\alpha}) \text{ is good} ~]
    \le \frac{1}{\iota (|\Field|-2)},
    \quad
    \forall j \in [\log_2 t].
\end{equation}

Fix a prefix $(\alpha_1, \ldots, \alpha_{j-1})$ such that $\bm{\xi}_{j-1}(\bm{\alpha})$ is good, and consider the set:
\[
    \mathbb{U}_j =
    \{ (1,\alpha_1) \kronecker \cdots \kronecker (1, \alpha_{j-1})
    \kronecker
    (1,\alpha) \kronecker (1,0) \kronecker \cdots \kronecker (1,0) \in \TensorSpace(t) ~|~ \alpha \in \Field \}.
\]
This is an affine space, since every vector $\bm{u}_\alpha \in \mathbb{U}_j$ can be written as:
\begin{align*}
\bm{u}_\alpha
=&
    (1,\alpha_1) \kronecker \cdots \kronecker (1, \alpha_{j-1})
    \kronecker
    (1,0) \kronecker (1,0) \kronecker \cdots \kronecker (1,0) 
\\
&
    + \alpha \cdot (1,\alpha_1) \kronecker \cdots \kronecker (1, \alpha_{j-1})
    \kronecker
    (0,1) \kronecker (1,0) \kronecker \cdots \kronecker (1,0).
\end{align*}
In particular, for $\alpha = 0$, we have $\bm{u}_0 = \bm{\xi}_{j-1}(\bm{\alpha})$, and for $\alpha = \alpha_j$, we obtain $\bm{u}_{\alpha_j} = \bm{\xi}_j(\bm{\alpha})$.

When $\bm{\xi}_{j-1}(\bm{\alpha})$ is good,
it holds that:
\[
    \Delta(\bm{u}_0,\CodeSpace) \ge (\epsilon - \iota \cdot (j-1))n.
\]
Applying Lemma~\ref{lemma:affinedistance} to the affine space $\mathbb{U}_j$, we conclude that:
\[
    \Pr_{\alpha \in \Field}[\Delta(\bm{u}_\alpha, \CodeSpace) < (\epsilon- \iota j)n] \le \frac{1}{\iota|\Field|}.
\]

Hence, conditioned on a good prefix $(\alpha_1, \ldots, \alpha_{j-1})$, there are at most $1/\iota$ choices of $\alpha_j \in \Field$ for which $\bm{\xi}_j(\bm{\alpha})$ is bad. Since $\alpha_j$ is sampled uniformly from $\Field \setminus \{0, -1\}$, which has cardinality $|\Field| - 2$, we obtain:
\[
    \Pr_{\bm{\alpha}}[~\bm{\xi}_{j}(\bm{\alpha}) \text{ is bad} \mid \bm{\xi}_{j-1}(\bm{\alpha}) \text{ is good} ~]
    \le \frac{1}{\iota (|\Field|-2)},
\]
This establishes Eq.~\eqref{eq:conditionalprobability}, and in turn implies Eq.~\eqref{eq:serialprobability}, thereby completing the proof of Eq.~\eqref{eq:projectrho}.
\end{proof}
